% Results of the clustering -- placeholders for figures, tables and interpretation.
% Replace each \figplaceholder by \includegraphics and each \hint by your own text.

\subsection{Results}                     % what the clustering found
\label{sec:clust_results}                % \secref{sec:clust_results}

\subsubsection{Agglomerative clustering}  % dendrogram, elbow, projections
\label{sec:clust_results_agglo}

\begin{figure}[htbp]                     % dendrogram + elbow curve
  \centering
  \figplaceholder[5cm]{Dendrogram of the Ward clustering (last 30 merges), with the cut line}
  \par\medskip
  \figplaceholder[4cm]{Average intra-cluster distance $D_{\mathrm{avg}}(k)$ and its relative reduction}
  \caption{Agglomerative clustering of the latent codes. Top: dendrogram, cut at $k^{*}$
           clusters (dashed line). Bottom: average intra-cluster distance as a function of the
           number of clusters (Eq.~\eqref{eq:davg}) and its relative reduction
           (Eq.~\eqref{eq:kstar}).}
  \label{fig:clust_agglo}
\end{figure}

\hint{Interpret: how many clusters are suggested, how clear is the largest drop compared with the
next candidates, and what are the cluster sizes? Does the dendrogram show well-separated
branches, or merges at similar heights?}

\begin{figure}[htbp]                     % the three projections, coloured by Ward cluster
  \centering
  \figplaceholder[4.5cm]{Kernel PCA \quad|\quad t-SNE \quad|\quad UMAP, coloured by Ward cluster}
  \caption{Latent codes projected with kernel PCA, t-SNE (perplexity \TODO{value}) and UMAP
           (\texttt{n\_neighbors} \TODO{value}), coloured by agglomerative cluster.}
  \label{fig:clust_agglo_proj}
\end{figure}

\hint{Interpret: are the clusters separated in all three projections, only in some, or in none?
Is the picture stable over the parameter sweep?}

\subsubsection{K-means and kernel K-means}  % choice of K, agreement between methods
\label{sec:clust_results_kmeans}

\begin{figure}[htbp]                     % elbow + silhouette for both methods
  \centering
  \figplaceholder[4cm]{K-means and kernel K-means: inertia (elbow) and mean silhouette for $K = 2, \dots, 10$}
  \caption{Choice of the number of clusters for K-means and kernel K-means: inertia
           (Eq.~\eqref{eq:kmeans}) and mean silhouette (Eq.~\eqref{eq:silhouette}).}
  \label{fig:clust_kmeans_k}
\end{figure}

\begin{table}[htbp]                      % agreement between the partitions
  \centering
  \small
  \caption{Agreement between partitions, measured with the adjusted Rand index (1 = same
           partition, 0 = chance level).}
  \label{tab:clust_ari}
  \begin{tabular}{llc}
    \toprule
    Comparison                                   & Latent codes          & ARI \\
    \midrule
    Agglomerative ($K=4$) vs K-means ($K=3$)     & unregistered          & \TODO{?} \\
    Agglomerative ($K=4$) vs K-means ($K=3$)     & registered            & \TODO{?} \\
    K-means: unregistered vs registered          & --                    & \TODO{?} \\
    Agglomerative: unregistered vs registered    & --                    & \TODO{?} \\
    \bottomrule
  \end{tabular}
\end{table}

\hint{Interpret: which $K$ do the elbow and the silhouette suggest, and do they agree? How much do
the two methods, and the two sets of codes, agree (Table~\ref{tab:clust_ari})? Report the robustness checks of
\secref{sec:clust_kmeans}: seed stability (mean ARI), sensitivity to $\gamma$, and the silhouette
of the real data compared with the no-cluster Gaussian baseline.}

\subsubsection{Morphological characterization of the clusters}  % what distinguishes the groups
\label{sec:clust_results_morph}

For every tree, the six descriptors of \secref{sec:morphometrics} are averaged over its branches
and complemented with two whole-tree counts, the number of branches and the deepest generation.
Comparing these eight values between clusters shows what distinguishes the groups, and colouring
the projections by each value shows whether the latent space is organized along these
anatomical properties.

\begin{figure}[htbp]                     % projections coloured by metric
  \centering
  \figplaceholder[5cm]{Projections coloured by each of the 8 metrics}
  \caption{UMAP projection of the latent codes coloured by the mean value of each
           morphological descriptor and by the two tree counts.}
  \label{fig:clust_metric_proj}
\end{figure}

\begin{figure}[htbp]                     % box plots per cluster
  \centering
  \figplaceholder[5cm]{Box plots of the 8 metrics per cluster}
  \caption{Distribution of the eight metrics in each cluster.}
  \label{fig:clust_boxplots}
\end{figure}

Tables~\ref{tab:unreg-kmeans} to~\ref{tab:reg-agglo} give, for each method and each set of
latent codes, the distribution of the eight metrics within every cluster: mean, standard deviation
and variance, median and interquartile range. The shading of each cell shows the standardized
difference between the cluster mean and the mean over all 419 trees,
\begin{equation}
  d = \frac{\bar{x}_{\text{cluster}} - \bar{x}_{\text{all}}}{s_{\text{all}}} ,
  \label{eq:cluster_d}                   % effect size used for the shading
\end{equation}
where $s_{\text{all}}$ is the standard deviation over all trees: red cells lie above the overall
mean, blue cells below, and darker cells further from it. Expressing the difference in units of
the overall spread makes metrics with very different scales comparable, from tortuosity values
around 0.6 to branch counts above 200. Whether a metric differs significantly between the
clusters can be tested with a Kruskal--Wallis test \cite{kruskal1952}, which compares the
distributions without assuming normality. \TODO{Add the Kruskal--Wallis $p$-values per metric,
e.g.\ as one sentence per table or as an extra column.}

% Per-cluster statistics of the 8 morphological metrics (tables from the clustering notebooks).
% \clustertablesetup (setup/helpers.tex) applies the table spacing locally; \topgap and \botgap
% (same file) are invisible struts that add space above the mean and below the SD line.

\begin{table}[htbp]
\centering
\clustertablesetup
\caption{Morphological metrics per cluster: K-means ($K=3$), latent codes of the unregistered trees.}
\label{tab:unreg-kmeans}
\resizebox{\linewidth}{!}{%
\begin{tabular}{lrrr|r}
\toprule
\topgap\textbf{Metric} & \textcolor{dotA}{\textbullet}\,\textbf{Cluster 0 (n=140)} & \textcolor{dotB}{\textbullet}\,\textbf{Cluster 1 (n=107)} & \textcolor{dotC}{\textbullet}\,\textbf{Cluster 2 (n=172)} & \textbf{All}\botgap \\
\midrule
\topgap\textbf{Divergence} & \cellcolor{loBlue!2}\topgap\textbf{0.3799} & \cellcolor{loBlue!14}\topgap\textbf{0.3645} & \cellcolor{hiRed!10}\topgap\textbf{0.3955} & \topgap\textbf{0.3824} \\
 & \cellcolor{loBlue!2}\stat SD 0.05221 $\cdot$ $\sigma^2$ 0.002726\botgap & \cellcolor{loBlue!14}\stat SD 0.05863 $\cdot$ $\sigma^2$ 0.003437\botgap & \cellcolor{hiRed!10}\stat SD 0.0451 $\cdot$ $\sigma^2$ 0.002034\botgap & \stat SD 0.05258 $\cdot$ $\sigma^2$ 0.002764\botgap \\
\arrayrulecolor{black!20}\hline\arrayrulecolor{black}
\topgap\textbf{Length} & \cellcolor{loBlue!3}\topgap\textbf{52.02} & \cellcolor{loBlue!19}\topgap\textbf{48.82} & \cellcolor{hiRed!14}\topgap\textbf{55.39} & \topgap\textbf{52.58} \\
 & \cellcolor{loBlue!3}\stat SD 7.471 $\cdot$ $\sigma^2$ 55.82\botgap & \cellcolor{loBlue!19}\stat SD 9.631 $\cdot$ $\sigma^2$ 92.77\botgap & \cellcolor{hiRed!14}\stat SD 6.287 $\cdot$ $\sigma^2$ 39.53\botgap & \stat SD 8.077 $\cdot$ $\sigma^2$ 65.23\botgap \\
\arrayrulecolor{black!20}\hline\arrayrulecolor{black}
\topgap\textbf{Complexity} & \cellcolor{loBlue!6}\topgap\textbf{0.7748} & \cellcolor{loBlue!11}\topgap\textbf{0.7643} & \cellcolor{hiRed!12}\topgap\textbf{0.811} & \topgap\textbf{0.787} \\
 & \cellcolor{loBlue!6}\stat SD 0.08176 $\cdot$ $\sigma^2$ 0.006685\botgap & \cellcolor{loBlue!11}\stat SD 0.1027 $\cdot$ $\sigma^2$ 0.01055\botgap & \cellcolor{hiRed!12}\stat SD 0.06113 $\cdot$ $\sigma^2$ 0.003737\botgap & \stat SD 0.08274 $\cdot$ $\sigma^2$ 0.006845\botgap \\
\arrayrulecolor{black!20}\hline\arrayrulecolor{black}
\topgap\textbf{Tortuosity} & \cellcolor{hiRed!11}\topgap\textbf{0.6253} & \cellcolor{loBlue!31}\topgap\textbf{0.6029} & \cellcolor{hiRed!10}\topgap\textbf{0.6247} & \topgap\textbf{0.6193} \\
 & \cellcolor{hiRed!11}\stat SD 0.01658 $\cdot$ $\sigma^2$ $2.75{\times}10^{-4}$\botgap & \cellcolor{loBlue!31}\stat SD 0.02305 $\cdot$ $\sigma^2$ $5.31{\times}10^{-4}$\botgap & \cellcolor{hiRed!10}\stat SD 0.01793 $\cdot$ $\sigma^2$ $3.21{\times}10^{-4}$\botgap & \stat SD 0.02123 $\cdot$ $\sigma^2$ $4.51{\times}10^{-4}$\botgap \\
\arrayrulecolor{black!20}\hline\arrayrulecolor{black}
\topgap\textbf{Stenosis} & \cellcolor{loBlue!7}\topgap\textbf{0.3748} & \cellcolor{hiRed!10}\topgap\textbf{0.3882} & \cellcolor{loBlue!1}\topgap\textbf{0.3799} & \topgap\textbf{0.3803} \\
 & \cellcolor{loBlue!7}\stat SD 0.03389 $\cdot$ $\sigma^2$ 0.001149\botgap & \cellcolor{hiRed!10}\stat SD 0.03415 $\cdot$ $\sigma^2$ 0.001166\botgap & \cellcolor{loBlue!1}\stat SD 0.0248 $\cdot$ $\sigma^2$ $6.15{\times}10^{-4}$\botgap & \stat SD 0.03091 $\cdot$ $\sigma^2$ $9.55{\times}10^{-4}$\botgap \\
\arrayrulecolor{black!20}\hline\arrayrulecolor{black}
\topgap\textbf{Ectasia} & \cellcolor{hiRed!1}\topgap\textbf{0.754} & \cellcolor{loBlue!28}\topgap\textbf{0.7433} & \cellcolor{hiRed!16}\topgap\textbf{0.7597} & \topgap\textbf{0.7536} \\
 & \cellcolor{hiRed!1}\stat SD 0.01499 $\cdot$ $\sigma^2$ $2.25{\times}10^{-4}$\botgap & \cellcolor{loBlue!28}\stat SD 0.01407 $\cdot$ $\sigma^2$ $1.98{\times}10^{-4}$\botgap & \cellcolor{hiRed!16}\stat SD 0.01118 $\cdot$ $\sigma^2$ $1.25{\times}10^{-4}$\botgap & \stat SD 0.01478 $\cdot$ $\sigma^2$ $2.18{\times}10^{-4}$\botgap \\
\arrayrulecolor{black!20}\hline\arrayrulecolor{black}
\topgap\textbf{Number of branches} & \cellcolor{loBlue!9}\topgap\textbf{217} & \topgap\textbf{234.6} & \cellcolor{hiRed!7}\topgap\textbf{250} & \topgap\textbf{235} \\
 & \cellcolor{loBlue!9}\stat SD 78.65 $\cdot$ $\sigma^2$ 6186\botgap & \stat SD 101.7 $\cdot$ $\sigma^2$ 10330\botgap & \cellcolor{hiRed!7}\stat SD 69.41 $\cdot$ $\sigma^2$ 4817\botgap & \stat SD 82.76 $\cdot$ $\sigma^2$ 6850\botgap \\
\arrayrulecolor{black!20}\hline\arrayrulecolor{black}
\topgap\textbf{Deepest generation} & \cellcolor{loBlue!6}\topgap\textbf{10.99} & \cellcolor{hiRed!2}\topgap\textbf{11.3} & \cellcolor{hiRed!4}\topgap\textbf{11.36} & \topgap\textbf{11.22} \\
 & \cellcolor{loBlue!6}\stat SD 1.476 $\cdot$ $\sigma^2$ 2.18\botgap & \cellcolor{hiRed!2}\stat SD 1.829 $\cdot$ $\sigma^2$ 3.344\botgap & \cellcolor{hiRed!4}\stat SD 1.279 $\cdot$ $\sigma^2$ 1.635\botgap & \stat SD 1.506 $\cdot$ $\sigma^2$ 2.269\botgap \\
\bottomrule
\end{tabular}}
\par\smallskip\footnotesize Each cell: mean (bold), standard deviation and variance. Shading: standardized difference $d=(\bar{x}_{\text{cluster}}-\bar{x}_{\text{all}})/s_{\text{all}}$, clipped at $|d|=1$; red = cluster mean above the overall mean, blue = below, darker = larger $|d|$.
\end{table}

\begin{table}[htbp]
\centering
\clustertablesetup
\caption{Morphological metrics per cluster: agglomerative clustering ($K=4$), latent codes of the unregistered trees.}
\label{tab:unreg-agglo}
\resizebox{\linewidth}{!}{%
\begin{tabular}{lrrrr|r}
\toprule
\topgap\textbf{Metric} & \textcolor{dotA}{\textbullet}\,\textbf{Cluster 0 (n=187)} & \textcolor{dotB}{\textbullet}\,\textbf{Cluster 1 (n=112)} & \textcolor{dotC}{\textbullet}\,\textbf{Cluster 2 (n=68)} & \textcolor{dotD}{\textbullet}\,\textbf{Cluster 3 (n=52)} & \textbf{All}\botgap \\
\midrule
\topgap\textbf{Divergence} & \topgap\textbf{0.3824} & \cellcolor{hiRed!15}\topgap\textbf{0.4018} & \cellcolor{loBlue!19}\topgap\textbf{0.3579} & \cellcolor{loBlue!8}\topgap\textbf{0.3724} & \topgap\textbf{0.3824} \\
 & \stat SD 0.05107 $\cdot$ $\sigma^2$ 0.002608\botgap & \cellcolor{hiRed!15}\stat SD 0.03565 $\cdot$ $\sigma^2$ 0.001271\botgap & \cellcolor{loBlue!19}\stat SD 0.06208 $\cdot$ $\sigma^2$ 0.003854\botgap & \cellcolor{loBlue!8}\stat SD 0.05986 $\cdot$ $\sigma^2$ 0.003583\botgap & \stat SD 0.05258 $\cdot$ $\sigma^2$ 0.002764\botgap \\
\arrayrulecolor{black!20}\hline\arrayrulecolor{black}
\topgap\textbf{Length} & \cellcolor{hiRed!1}\topgap\textbf{52.82} & \cellcolor{hiRed!18}\topgap\textbf{56.15} & \cellcolor{loBlue!21}\topgap\textbf{48.37} & \cellcolor{loBlue!15}\topgap\textbf{49.58} & \topgap\textbf{52.58} \\
 & \cellcolor{hiRed!1}\stat SD 6.997 $\cdot$ $\sigma^2$ 48.95\botgap & \cellcolor{hiRed!18}\stat SD 5.685 $\cdot$ $\sigma^2$ 32.32\botgap & \cellcolor{loBlue!21}\stat SD 10.15 $\cdot$ $\sigma^2$ 103.1\botgap & \cellcolor{loBlue!15}\stat SD 9.452 $\cdot$ $\sigma^2$ 89.34\botgap & \stat SD 8.077 $\cdot$ $\sigma^2$ 65.23\botgap \\
\arrayrulecolor{black!20}\hline\arrayrulecolor{black}
\topgap\textbf{Complexity} & \cellcolor{hiRed!1}\topgap\textbf{0.79} & \cellcolor{hiRed!15}\topgap\textbf{0.8175} & \cellcolor{loBlue!16}\topgap\textbf{0.7535} & \cellcolor{loBlue!16}\topgap\textbf{0.754} & \topgap\textbf{0.787} \\
 & \cellcolor{hiRed!1}\stat SD 0.07292 $\cdot$ $\sigma^2$ 0.005317\botgap & \cellcolor{hiRed!15}\stat SD 0.05175 $\cdot$ $\sigma^2$ 0.002678\botgap & \cellcolor{loBlue!16}\stat SD 0.1084 $\cdot$ $\sigma^2$ 0.01175\botgap & \cellcolor{loBlue!16}\stat SD 0.1037 $\cdot$ $\sigma^2$ 0.01076\botgap & \stat SD 0.08274 $\cdot$ $\sigma^2$ 0.006845\botgap \\
\arrayrulecolor{black!20}\hline\arrayrulecolor{black}
\topgap\textbf{Tortuosity} & \cellcolor{loBlue!2}\topgap\textbf{0.6185} & \cellcolor{hiRed!14}\topgap\textbf{0.6265} & \cellcolor{loBlue!23}\topgap\textbf{0.607} & \cellcolor{hiRed!7}\topgap\textbf{0.623} & \topgap\textbf{0.6193} \\
 & \cellcolor{loBlue!2}\stat SD 0.02183 $\cdot$ $\sigma^2$ $4.76{\times}10^{-4}$\botgap & \cellcolor{hiRed!14}\stat SD 0.01584 $\cdot$ $\sigma^2$ $2.51{\times}10^{-4}$\botgap & \cellcolor{loBlue!23}\stat SD 0.02483 $\cdot$ $\sigma^2$ $6.17{\times}10^{-4}$\botgap & \cellcolor{hiRed!7}\stat SD 0.01636 $\cdot$ $\sigma^2$ $2.68{\times}10^{-4}$\botgap & \stat SD 0.02123 $\cdot$ $\sigma^2$ $4.51{\times}10^{-4}$\botgap \\
\arrayrulecolor{black!20}\hline\arrayrulecolor{black}
\topgap\textbf{Stenosis} & \cellcolor{hiRed!2}\topgap\textbf{0.3816} & \cellcolor{loBlue!1}\topgap\textbf{0.3796} & \cellcolor{hiRed!3}\topgap\textbf{0.3829} & \cellcolor{loBlue!8}\topgap\textbf{0.3739} & \topgap\textbf{0.3803} \\
 & \cellcolor{hiRed!2}\stat SD 0.0283 $\cdot$ $\sigma^2$ $8.01{\times}10^{-4}$\botgap & \cellcolor{loBlue!1}\stat SD 0.02333 $\cdot$ $\sigma^2$ $5.44{\times}10^{-4}$\botgap & \cellcolor{hiRed!3}\stat SD 0.04191 $\cdot$ $\sigma^2$ 0.001756\botgap & \cellcolor{loBlue!8}\stat SD 0.03687 $\cdot$ $\sigma^2$ 0.001359\botgap & \stat SD 0.03091 $\cdot$ $\sigma^2$ $9.55{\times}10^{-4}$\botgap \\
\arrayrulecolor{black!20}\hline\arrayrulecolor{black}
\topgap\textbf{Ectasia} & \cellcolor{loBlue!4}\topgap\textbf{0.7522} & \cellcolor{hiRed!19}\topgap\textbf{0.7605} & \cellcolor{loBlue!15}\topgap\textbf{0.7481} & \cellcolor{loBlue!7}\topgap\textbf{0.7509} & \topgap\textbf{0.7536} \\
 & \cellcolor{loBlue!4}\stat SD 0.01548 $\cdot$ $\sigma^2$ $2.4{\times}10^{-4}$\botgap & \cellcolor{hiRed!19}\stat SD 0.009722 $\cdot$ $\sigma^2$ $9.45{\times}10^{-5}$\botgap & \cellcolor{loBlue!15}\stat SD 0.01402 $\cdot$ $\sigma^2$ $1.96{\times}10^{-4}$\botgap & \cellcolor{loBlue!7}\stat SD 0.01706 $\cdot$ $\sigma^2$ $2.91{\times}10^{-4}$\botgap & \stat SD 0.01478 $\cdot$ $\sigma^2$ $2.18{\times}10^{-4}$\botgap \\
\arrayrulecolor{black!20}\hline\arrayrulecolor{black}
\topgap\textbf{Number of branches} & \topgap\textbf{236} & \cellcolor{hiRed!9}\topgap\textbf{253.1} & \cellcolor{loBlue!6}\topgap\textbf{222.9} & \cellcolor{loBlue!13}\topgap\textbf{208.4} & \topgap\textbf{235} \\
 & \stat SD 81.53 $\cdot$ $\sigma^2$ 6648\botgap & \cellcolor{hiRed!9}\stat SD 58.7 $\cdot$ $\sigma^2$ 3445\botgap & \cellcolor{loBlue!6}\stat SD 104.7 $\cdot$ $\sigma^2$ 10960\botgap & \cellcolor{loBlue!13}\stat SD 91.41 $\cdot$ $\sigma^2$ 8356\botgap & \stat SD 82.76 $\cdot$ $\sigma^2$ 6850\botgap \\
\arrayrulecolor{black!20}\hline\arrayrulecolor{black}
\topgap\textbf{Deepest generation} & \cellcolor{hiRed!2}\topgap\textbf{11.29} & \cellcolor{hiRed!6}\topgap\textbf{11.44} & \cellcolor{hiRed!1}\topgap\textbf{11.25} & \cellcolor{loBlue!20}\topgap\textbf{10.48} & \topgap\textbf{11.22} \\
 & \cellcolor{hiRed!2}\stat SD 1.489 $\cdot$ $\sigma^2$ 2.217\botgap & \cellcolor{hiRed!6}\stat SD 1.129 $\cdot$ $\sigma^2$ 1.275\botgap & \cellcolor{hiRed!1}\stat SD 1.872 $\cdot$ $\sigma^2$ 3.504\botgap & \cellcolor{loBlue!20}\stat SD 1.565 $\cdot$ $\sigma^2$ 2.451\botgap & \stat SD 1.506 $\cdot$ $\sigma^2$ 2.269\botgap \\
\bottomrule
\end{tabular}}
\par\smallskip\footnotesize Each cell: mean (bold), standard deviation and variance. Shading: standardized difference $d=(\bar{x}_{\text{cluster}}-\bar{x}_{\text{all}})/s_{\text{all}}$, clipped at $|d|=1$; red = cluster mean above the overall mean, blue = below, darker = larger $|d|$.
\end{table}

\begin{table}[htbp]
\centering
\clustertablesetup
\caption{Morphological metrics per cluster: K-means ($K=3$), latent codes of the registered trees.}
\label{tab:reg-kmeans}
\resizebox{\linewidth}{!}{%
\begin{tabular}{lrrr|r}
\toprule
\topgap\textbf{Metric} & \textcolor{dotA}{\textbullet}\,\textbf{Cluster 0 (n=151)} & \textcolor{dotB}{\textbullet}\,\textbf{Cluster 1 (n=200)} & \textcolor{dotC}{\textbullet}\,\textbf{Cluster 2 (n=68)} & \textbf{All}\botgap \\
\midrule
\topgap\textbf{Divergence} & \cellcolor{hiRed!16}\topgap\textbf{0.4032} & \cellcolor{loBlue!11}\topgap\textbf{0.3676} & \cellcolor{loBlue!2}\topgap\textbf{0.3794} & \topgap\textbf{0.3824} \\
 & \cellcolor{hiRed!16}\stat SD 0.0438 $\cdot$ $\sigma^2$ 0.001918\botgap & \cellcolor{loBlue!11}\stat SD 0.0522 $\cdot$ $\sigma^2$ 0.002725\botgap & \cellcolor{loBlue!2}\stat SD 0.05677 $\cdot$ $\sigma^2$ 0.003222\botgap & \stat SD 0.05258 $\cdot$ $\sigma^2$ 0.002764\botgap \\
\arrayrulecolor{black!20}\hline\arrayrulecolor{black}
\topgap\textbf{Length} & \cellcolor{hiRed!18}\topgap\textbf{56.19} & \cellcolor{loBlue!14}\topgap\textbf{49.78} & \cellcolor{hiRed!1}\topgap\textbf{52.85} & \topgap\textbf{52.58} \\
 & \cellcolor{hiRed!18}\stat SD 6.06 $\cdot$ $\sigma^2$ 36.73\botgap & \cellcolor{loBlue!14}\stat SD 8.307 $\cdot$ $\sigma^2$ 69\botgap & \cellcolor{hiRed!1}\stat SD 8.182 $\cdot$ $\sigma^2$ 66.94\botgap & \stat SD 8.077 $\cdot$ $\sigma^2$ 65.23\botgap \\
\arrayrulecolor{black!20}\hline\arrayrulecolor{black}
\topgap\textbf{Complexity} & \cellcolor{hiRed!16}\topgap\textbf{0.821} & \cellcolor{loBlue!12}\topgap\textbf{0.7613} & \topgap\textbf{0.7869} & \topgap\textbf{0.787} \\
 & \cellcolor{hiRed!16}\stat SD 0.06085 $\cdot$ $\sigma^2$ 0.003703\botgap & \cellcolor{loBlue!12}\stat SD 0.08631 $\cdot$ $\sigma^2$ 0.00745\botgap & \stat SD 0.08787 $\cdot$ $\sigma^2$ 0.007722\botgap & \stat SD 0.08274 $\cdot$ $\sigma^2$ 0.006845\botgap \\
\arrayrulecolor{black!20}\hline\arrayrulecolor{black}
\topgap\textbf{Tortuosity} & \cellcolor{hiRed!5}\topgap\textbf{0.622} & \cellcolor{hiRed!1}\topgap\textbf{0.6197} & \cellcolor{loBlue!13}\topgap\textbf{0.6124} & \topgap\textbf{0.6193} \\
 & \cellcolor{hiRed!5}\stat SD 0.01903 $\cdot$ $\sigma^2$ $3.62{\times}10^{-4}$\botgap & \cellcolor{hiRed!1}\stat SD 0.02246 $\cdot$ $\sigma^2$ $5.04{\times}10^{-4}$\botgap & \cellcolor{loBlue!13}\stat SD 0.02094 $\cdot$ $\sigma^2$ $4.39{\times}10^{-4}$\botgap & \stat SD 0.02123 $\cdot$ $\sigma^2$ $4.51{\times}10^{-4}$\botgap \\
\arrayrulecolor{black!20}\hline\arrayrulecolor{black}
\topgap\textbf{Stenosis} & \topgap\textbf{0.38} & \cellcolor{loBlue!4}\topgap\textbf{0.3772} & \cellcolor{hiRed!13}\topgap\textbf{0.3902} & \topgap\textbf{0.3803} \\
 & \stat SD 0.02613 $\cdot$ $\sigma^2$ $6.83{\times}10^{-4}$\botgap & \cellcolor{loBlue!4}\stat SD 0.03422 $\cdot$ $\sigma^2$ 0.001171\botgap & \cellcolor{hiRed!13}\stat SD 0.02871 $\cdot$ $\sigma^2$ $8.24{\times}10^{-4}$\botgap & \stat SD 0.03091 $\cdot$ $\sigma^2$ $9.55{\times}10^{-4}$\botgap \\
\arrayrulecolor{black!20}\hline\arrayrulecolor{black}
\topgap\textbf{Ectasia} & \cellcolor{hiRed!14}\topgap\textbf{0.7589} & \cellcolor{loBlue!7}\topgap\textbf{0.751} & \cellcolor{loBlue!11}\topgap\textbf{0.7494} & \topgap\textbf{0.7536} \\
 & \cellcolor{hiRed!14}\stat SD 0.01328 $\cdot$ $\sigma^2$ $1.76{\times}10^{-4}$\botgap & \cellcolor{loBlue!7}\stat SD 0.01537 $\cdot$ $\sigma^2$ $2.36{\times}10^{-4}$\botgap & \cellcolor{loBlue!11}\stat SD 0.01278 $\cdot$ $\sigma^2$ $1.63{\times}10^{-4}$\botgap & \stat SD 0.01478 $\cdot$ $\sigma^2$ $2.18{\times}10^{-4}$\botgap \\
\arrayrulecolor{black!20}\hline\arrayrulecolor{black}
\topgap\textbf{Number of branches} & \cellcolor{hiRed!14}\topgap\textbf{263.8} & \cellcolor{loBlue!11}\topgap\textbf{212.9} & \cellcolor{hiRed!1}\topgap\textbf{236.4} & \topgap\textbf{235} \\
 & \cellcolor{hiRed!14}\stat SD 74.84 $\cdot$ $\sigma^2$ 5601\botgap & \cellcolor{loBlue!11}\stat SD 77.67 $\cdot$ $\sigma^2$ 6032\botgap & \cellcolor{hiRed!1}\stat SD 94.54 $\cdot$ $\sigma^2$ 8938\botgap & \stat SD 82.76 $\cdot$ $\sigma^2$ 6850\botgap \\
\arrayrulecolor{black!20}\hline\arrayrulecolor{black}
\topgap\textbf{Deepest generation} & \cellcolor{hiRed!8}\topgap\textbf{11.53} & \cellcolor{loBlue!7}\topgap\textbf{10.96} & \cellcolor{hiRed!2}\topgap\textbf{11.31} & \topgap\textbf{11.22} \\
 & \cellcolor{hiRed!8}\stat SD 1.356 $\cdot$ $\sigma^2$ 1.837\botgap & \cellcolor{loBlue!7}\stat SD 1.49 $\cdot$ $\sigma^2$ 2.219\botgap & \cellcolor{hiRed!2}\stat SD 1.739 $\cdot$ $\sigma^2$ 3.023\botgap & \stat SD 1.506 $\cdot$ $\sigma^2$ 2.269\botgap \\
\bottomrule
\end{tabular}}
\par\smallskip\footnotesize Each cell: mean (bold), standard deviation and variance. Shading: standardized difference $d=(\bar{x}_{\text{cluster}}-\bar{x}_{\text{all}})/s_{\text{all}}$, clipped at $|d|=1$; red = cluster mean above the overall mean, blue = below, darker = larger $|d|$.
\end{table}

\begin{table}[htbp]
\centering
\clustertablesetup
\caption{Morphological metrics per cluster: agglomerative clustering ($K=4$), latent codes of the registered trees.}
\label{tab:reg-agglo}
\resizebox{\linewidth}{!}{%
\begin{tabular}{lrrrr|r}
\toprule
\topgap\textbf{Metric} & \textcolor{dotA}{\textbullet}\,\textbf{Cluster 0 (n=325)} & \textcolor{dotB}{\textbullet}\,\textbf{Cluster 1 (n=74)} & \textcolor{dotC}{\textbullet}\,\textbf{Cluster 2 (n=13)} & \textcolor{dotD}{\textbullet}\,\textbf{Cluster 3 (n=7)} & \textbf{All}\botgap \\
\midrule
\topgap\textbf{Divergence} & \cellcolor{loBlue!3}\topgap\textbf{0.3789} & \cellcolor{hiRed!8}\topgap\textbf{0.3931} & \cellcolor{hiRed!2}\topgap\textbf{0.3854} & \cellcolor{hiRed!32}\topgap\textbf{0.4242} & \topgap\textbf{0.3824} \\
 & \cellcolor{loBlue!3}\stat SD 0.05366 $\cdot$ $\sigma^2$ 0.002879\botgap & \cellcolor{hiRed!8}\stat SD 0.04695 $\cdot$ $\sigma^2$ 0.002205\botgap & \cellcolor{hiRed!2}\stat SD 0.04919 $\cdot$ $\sigma^2$ 0.00242\botgap & \cellcolor{hiRed!32}\stat SD 0.03629 $\cdot$ $\sigma^2$ 0.001317\botgap & \stat SD 0.05258 $\cdot$ $\sigma^2$ 0.002764\botgap \\
\arrayrulecolor{black!20}\hline\arrayrulecolor{black}
\topgap\textbf{Length} & \cellcolor{loBlue!2}\topgap\textbf{52.27} & \cellcolor{hiRed!4}\topgap\textbf{53.34} & \cellcolor{loBlue!1}\topgap\textbf{52.42} & \cellcolor{hiRed!35}\topgap\textbf{59.61} & \topgap\textbf{52.58} \\
 & \cellcolor{loBlue!2}\stat SD 8.174 $\cdot$ $\sigma^2$ 66.82\botgap & \cellcolor{hiRed!4}\stat SD 6.939 $\cdot$ $\sigma^2$ 48.15\botgap & \cellcolor{loBlue!1}\stat SD 11.04 $\cdot$ $\sigma^2$ 121.9\botgap & \cellcolor{hiRed!35}\stat SD 6.027 $\cdot$ $\sigma^2$ 36.33\botgap & \stat SD 8.077 $\cdot$ $\sigma^2$ 65.23\botgap \\
\arrayrulecolor{black!20}\hline\arrayrulecolor{black}
\topgap\textbf{Complexity} & \cellcolor{loBlue!2}\topgap\textbf{0.7822} & \cellcolor{hiRed!8}\topgap\textbf{0.8042} & \cellcolor{loBlue!8}\topgap\textbf{0.7707} & \cellcolor{hiRed!35}\topgap\textbf{0.8591} & \topgap\textbf{0.787} \\
 & \cellcolor{loBlue!2}\stat SD 0.08453 $\cdot$ $\sigma^2$ 0.007145\botgap & \cellcolor{hiRed!8}\stat SD 0.06708 $\cdot$ $\sigma^2$ 0.004499\botgap & \cellcolor{loBlue!8}\stat SD 0.1059 $\cdot$ $\sigma^2$ 0.01121\botgap & \cellcolor{hiRed!35}\stat SD 0.04858 $\cdot$ $\sigma^2$ 0.00236\botgap & \stat SD 0.08274 $\cdot$ $\sigma^2$ 0.006845\botgap \\
\arrayrulecolor{black!20}\hline\arrayrulecolor{black}
\topgap\textbf{Tortuosity} & \topgap\textbf{0.6195} & \cellcolor{loBlue!3}\topgap\textbf{0.6177} & \cellcolor{hiRed!8}\topgap\textbf{0.6235} & \cellcolor{hiRed!3}\topgap\textbf{0.621} & \topgap\textbf{0.6193} \\
 & \stat SD 0.02139 $\cdot$ $\sigma^2$ $4.58{\times}10^{-4}$\botgap & \cellcolor{loBlue!3}\stat SD 0.02139 $\cdot$ $\sigma^2$ $4.57{\times}10^{-4}$\botgap & \cellcolor{hiRed!8}\stat SD 0.0197 $\cdot$ $\sigma^2$ $3.88{\times}10^{-4}$\botgap & \cellcolor{hiRed!3}\stat SD 0.01674 $\cdot$ $\sigma^2$ $2.8{\times}10^{-4}$\botgap & \stat SD 0.02123 $\cdot$ $\sigma^2$ $4.51{\times}10^{-4}$\botgap \\
\arrayrulecolor{black!20}\hline\arrayrulecolor{black}
\topgap\textbf{Stenosis} & \topgap\textbf{0.3802} & \cellcolor{loBlue!1}\topgap\textbf{0.3797} & \cellcolor{loBlue!2}\topgap\textbf{0.3789} & \cellcolor{hiRed!17}\topgap\textbf{0.3937} & \topgap\textbf{0.3803} \\
 & \stat SD 0.03151 $\cdot$ $\sigma^2$ $9.93{\times}10^{-4}$\botgap & \cellcolor{loBlue!1}\stat SD 0.02922 $\cdot$ $\sigma^2$ $8.54{\times}10^{-4}$\botgap & \cellcolor{loBlue!2}\stat SD 0.03077 $\cdot$ $\sigma^2$ $9.47{\times}10^{-4}$\botgap & \cellcolor{hiRed!17}\stat SD 0.02064 $\cdot$ $\sigma^2$ $4.26{\times}10^{-4}$\botgap & \stat SD 0.03091 $\cdot$ $\sigma^2$ $9.55{\times}10^{-4}$\botgap \\
\arrayrulecolor{black!20}\hline\arrayrulecolor{black}
\topgap\textbf{Ectasia} & \cellcolor{loBlue!1}\topgap\textbf{0.7532} & \cellcolor{hiRed!1}\topgap\textbf{0.7541} & \cellcolor{hiRed!4}\topgap\textbf{0.755} & \cellcolor{hiRed!31}\topgap\textbf{0.765} & \topgap\textbf{0.7536} \\
 & \cellcolor{loBlue!1}\stat SD 0.01518 $\cdot$ $\sigma^2$ $2.3{\times}10^{-4}$\botgap & \cellcolor{hiRed!1}\stat SD 0.01261 $\cdot$ $\sigma^2$ $1.59{\times}10^{-4}$\botgap & \cellcolor{hiRed!4}\stat SD 0.0177 $\cdot$ $\sigma^2$ $3.13{\times}10^{-4}$\botgap & \cellcolor{hiRed!31}\stat SD 0.006023 $\cdot$ $\sigma^2$ $3.63{\times}10^{-5}$\botgap & \stat SD 0.01478 $\cdot$ $\sigma^2$ $2.18{\times}10^{-4}$\botgap \\
\arrayrulecolor{black!20}\hline\arrayrulecolor{black}
\topgap\textbf{Number of branches} & \cellcolor{loBlue!2}\topgap\textbf{230.4} & \cellcolor{hiRed!8}\topgap\textbf{252} & \cellcolor{loBlue!7}\topgap\textbf{220.5} & \cellcolor{hiRed!31}\topgap\textbf{299.3} & \topgap\textbf{235} \\
 & \cellcolor{loBlue!2}\stat SD 84.18 $\cdot$ $\sigma^2$ 7087\botgap & \cellcolor{hiRed!8}\stat SD 71.89 $\cdot$ $\sigma^2$ 5167\botgap & \cellcolor{loBlue!7}\stat SD 83.75 $\cdot$ $\sigma^2$ 7013\botgap & \cellcolor{hiRed!31}\stat SD 87.31 $\cdot$ $\sigma^2$ 7624\botgap & \stat SD 82.76 $\cdot$ $\sigma^2$ 6850\botgap \\
\arrayrulecolor{black!20}\hline\arrayrulecolor{black}
\topgap\textbf{Deepest generation} & \cellcolor{loBlue!1}\topgap\textbf{11.19} & \cellcolor{hiRed!4}\topgap\textbf{11.38} & \cellcolor{loBlue!6}\topgap\textbf{11} & \cellcolor{hiRed!9}\topgap\textbf{11.57} & \topgap\textbf{11.22} \\
 & \cellcolor{loBlue!1}\stat SD 1.541 $\cdot$ $\sigma^2$ 2.375\botgap & \cellcolor{hiRed!4}\stat SD 1.372 $\cdot$ $\sigma^2$ 1.882\botgap & \cellcolor{loBlue!6}\stat SD 1.528 $\cdot$ $\sigma^2$ 2.333\botgap & \cellcolor{hiRed!9}\stat SD 1.272 $\cdot$ $\sigma^2$ 1.619\botgap & \stat SD 1.506 $\cdot$ $\sigma^2$ 2.269\botgap \\
\bottomrule
\end{tabular}}
\par\smallskip\footnotesize Each cell: mean (bold), standard deviation and variance. Shading: standardized difference $d=(\bar{x}_{\text{cluster}}-\bar{x}_{\text{all}})/s_{\text{all}}$, clipped at $|d|=1$; red = cluster mean above the overall mean, blue = below, darker = larger $|d|$.
\end{table}  % the four landscape tables

\hint{Interpret: which metrics separate the clusters most (darkest cells), and in which
direction? Can the clusters be described in anatomical terms, e.g.\ longer, more complex and
more branched trees against shorter, simpler ones? Compare unregistered and registered codes:
does registration change which properties drive the clusters? Note also the cluster sizes: with
registered codes, agglomerative clustering produces two very small clusters (13 and 7 trees)
next to one of 325 trees. Do the metric-coloured projections show gradients that follow the
clusters, or that cut across them?}